% © 2026 Ing. David Jaroš — CC BY-NC-ND 4.0
%
% This work is licensed under a Creative Commons Attribution-NonCommercial-NoDerivatives
% 4.0 International License (CC BY-NC-ND 4.0).
%
% License History: Earlier drafts (up to v0.3) were released under CC BY 4.0.
% From v0.4 onward, all material is released under CC BY-NC-ND 4.0 to protect
% the integrity of the theoretical work during ongoing academic development.

% =====================================================================
% TASK 8 — R Factor Two-Loop Analysis
% Status: [OPEN HARD PROBLEM] — B3 approach, coefficient c not derived
% Source: DERIVATION_INDEX.md §"Fine Structure Constant" R-factor row
% =====================================================================

\documentclass[11pt,a4paper]{article}
\usepackage{amsmath, amssymb, amsthm}
\usepackage{hyperref}
\usepackage{geometry}
\geometry{margin=2.5cm}

\newtheorem{theorem}{Theorem}
\newtheorem{proposition}{Proposition}
\newtheorem{conjecture}{Conjecture}
\theoremstyle{remark}
\newtheorem{remark}{Remark}

\newcommand{\Lzero}{\textbf{[L0]}}
\newcommand{\Lone}{\textbf{[L1]}}
\newcommand{\OPEN}{\textbf{[OPEN]}}

\title{\textbf{R Factor: Two-Loop Analysis on the $\psi$-Circle}}
\author{Ing. David Jaroš}
\date{March 2026 (v57)}

\begin{document}
\maketitle

\begin{abstract}
The correction factor $R \approx 1.114$ in the UBT fine-structure derivation
($B = B_{\mathrm{base}} \times R$) is an open hard problem.  The best
algebraic candidate is $R = 1 + \alpha(N_{\mathrm{eff}} + \pi + 1/4) \approx 1.1123$
(error $0.15\%$), found numerically (\texttt{tools/explore\_r\_factor.py}).  Neither
the coefficient $c = N_{\mathrm{eff}} + \pi + 1/4 \approx 15.39$ nor the formula
$R = 1 + c\alpha$ has been derived from $\mathcal{S}[\Theta]$.  This document
analyses the two-loop beta-function approach (Approach~B3) and documents the
result.
\end{abstract}

% ======================================================================
\section{Setup: The $R$ Factor Problem}
\label{sec:setup}
% ======================================================================

The UBT prediction for $\alpha^{-1}$ requires:
\begin{equation}
  \alpha = \frac{B_{\mathrm{base}} \times R}{n^{*\,2}},
  \qquad
  B_{\mathrm{base}} = N_{\mathrm{eff}}^{3/2} \approx 41.57,
  \qquad n^* = 137.
  \label{eq:alpha_formula}
\end{equation}
With $\alpha = 1/137.036$:
\begin{equation}
  R = \frac{\alpha \cdot n^{*\,2}}{B_{\mathrm{base}}}
    = \frac{137.036^2}{137^2 \cdot 41.57 \cdot 137.036}
    \approx 1.114.
  \label{eq:R_target}
\end{equation}

\noindent\textbf{Status of $R$:}
\begin{center}
\begin{tabular}{|l|l|}
\hline
$B_{\mathrm{base}} = N_{\mathrm{eff}}^{3/2}$ & Partially Proved [L1] \\
$n^* = 137$ (prime attractor) & Proved [L1] \\
$R \approx 1.114$ & \OPEN\ Hard Problem \\
\hline
\end{tabular}
\end{center}

% ======================================================================
\section{Current State: \texttt{explore\_r\_factor.py}}
\label{sec:current}
% ======================================================================

Running \texttt{tools/explore\_r\_factor.py} gives the closest algebraic
candidates (sorted by error):

\begin{center}
\begin{tabular}{lll}
\hline
\textbf{Candidate} & \textbf{Value} & \textbf{Error} \\
\hline
$1 + \alpha(N_{\mathrm{eff}} + \pi + 1/4)$ & 1.11235 & 0.148\% \\
$1 + \alpha N_{\mathrm{eff}}^{3/2}/e$ & 1.11162 & 0.213\% \\
$1 + \alpha(N_{\mathrm{eff}} + \pi)$ & 1.11052 & 0.312\% \\
$2^{1/6}$ & 1.12246 & 0.760\% \\
\hline
\end{tabular}
\end{center}

Best candidate: $R = 1 + \alpha(N_{\mathrm{eff}} + \pi + 1/4)$.
This is a \textbf{Numerical Observation} only — not derived from $\mathcal{S}[\Theta]$.

% ======================================================================
\section{Two-Loop Beta Function on the $\psi$-Circle (Approach B3)}
\label{sec:twoloop}
% ======================================================================

\subsection{One-Loop Structure}

At one loop, the renormalisation group equation on the $\psi$-circle is:
\begin{equation}
  \mu\frac{d\alpha}{d\mu} = \beta_1\,\alpha^2 + \mathcal{O}(\alpha^3),
  \qquad \beta_1 = -\frac{2N_{\mathrm{eff}}}{3\pi}.
  \label{eq:one_loop_beta}
\end{equation}
The one-loop running gives the standard QED-like correction:
$\alpha(M_Z^2) \approx 1/128$, consistent with the observed running of $\alpha$.

\subsection{Two-Loop Contribution}

At two loops, the beta function is:
\begin{equation}
  \mu\frac{d\alpha}{d\mu}
    = \beta_1\,\alpha^2 + \beta_2\,\alpha^3 + \mathcal{O}(\alpha^4),
  \label{eq:two_loop_beta}
\end{equation}
where $\beta_2$ arises from the two-loop diagram on the $\psi$-circle.

The two-loop correction to the effective coupling is:
\begin{equation}
  \alpha_{\mathrm{eff}}
    = \alpha_0\bigl(1 + c\,\alpha_0\bigr) + \mathcal{O}(\alpha_0^3),
  \label{eq:alpha_two_loop}
\end{equation}
where $c$ is the two-loop coefficient.

\subsection{Identifying $c$}

The best numerical candidate gives $c \approx N_{\mathrm{eff}} + \pi + 1/4 = 15.39$.

Decomposing this:
\begin{itemize}
  \item $N_{\mathrm{eff}} = 12$: counting effective bosonic modes on $\psi$-circle.
  \item $\pi$: candidate origin — volume factor from angular integration over
    the $\psi$-circle half-period ($\mathrm{vol}(S^1)/(2\cdot 1) = \pi$).
  \item $1/4$: candidate origin — Weyl spinor loop factor
    ($2$ real DOF / $4D$ spacetime = $1/2$ per DOF; $1/2 \times 1/2 = 1/4$
    for two fermionic lines in the two-loop diagram).
\end{itemize}

\begin{conjecture}[Two-loop coefficient, B3]
\label{conj:two_loop}
The two-loop coefficient on the $\psi$-circle is:
\begin{equation}
  c = N_{\mathrm{eff}} + \pi + \frac{1}{4},
  \label{eq:c_conjecture}
\end{equation}
where $\pi$ comes from half-period integration and $1/4$ from the Weyl
spinor loop factor.  This gives $R = 1 + c\alpha \approx 1.1123$.
\end{conjecture}

\textbf{Status of Conjecture~\ref{conj:two_loop}:}
\textbf{Numerical Observation only} — neither the $\pi$ nor the $1/4$
contribution has been derived from an explicit two-loop Feynman diagram
computation in $\mathcal{S}[\Theta]$ on the $\psi$-circle.

% ======================================================================
\section{Two-Scale Approach (v70): Dead End}
\label{sec:two_scale}
% ======================================================================

\subsection{Setup}

The two $\psi$-circle radii established in UBT are:
\begin{equation}
  R_\psi^{(\alpha)} = \frac{\hbar}{n_R\,m_e\,c} \approx 2.80\ \text{fm},
  \qquad
  R_\psi^{(\mathrm{EW})} = \frac{n_R\,\pi}{2g\,v_{\mathrm{obs}}} \approx 0.27\ \text{fm}.
  \label{eq:two_scale_radii}
\end{equation}
Their ratio is $R_\psi^{(\alpha)}/R_\psi^{(\mathrm{EW})} \approx 10.4$.  The
question: can this ratio generate $R \approx 1.114$?

\subsection{Key Computation}

Setting $\rho = R_\psi^{(\mathrm{EW})}/R_\psi^{(\alpha)} \approx 0.096$, one
observes the numerical coincidence:
\begin{equation}
  \rho \approx 0.096 \approx \alpha \cdot N_{\mathrm{eff}}
  \quad\text{(numerical observation only).}
\end{equation}
Attempting $R = 1 + \rho \cdot \alpha \cdot \beta$ requires
$\beta \approx 163$ — no algebraic value identified.
Alternatively, $R = \rho^\varepsilon$ gives $\varepsilon \approx 0.046$
— again no algebraic value.

\subsection{Verdict}

\begin{center}
\textbf{Dead End} — the ratio $R_\psi^{(\alpha)}/R_\psi^{(\mathrm{EW})}$ does
not generate $R \approx 1.114$ through any natural algebraic combination.
\end{center}

% ======================================================================
\section{Modular Correction Approach (v70): Dead End}
\label{sec:modular}
% ======================================================================

\subsection{Setup}

Near the self-dual point $\tau = i$ of the $\psi$-circle modular parameter,
the partition function $Z(\tau)$ receives corrections from the Eisenstein series
$E_2(\tau)$.  Setting $\tau = i + \varepsilon$ with $\varepsilon \ll 1$:
\begin{equation}
  Z(\tau) \approx Z(i)\!\left[1 + \frac{3\pi\,\varepsilon}{2} + \mathcal{O}(\varepsilon^2)\right].
\end{equation}

\subsection{Key Computation}

With $\varepsilon = 1/(2\pi \cdot 137)$ (one-loop shift at the fine-structure
scale):
\begin{equation}
  \text{correction} \approx 1 + \frac{3\pi}{2}\cdot\frac{1}{2\pi\cdot 137}
  \approx 1.0055.
\end{equation}
This is a factor $\approx 20$ too small compared to $R - 1 \approx 0.114$.
No value of $\varepsilon$ derived from UBT parameters reproduces $R$.

\subsection{Verdict}

\begin{center}
\textbf{Dead End} — modular correction near $\tau = i$ is of order $\alpha \sim 1/137$,
yielding only $R \approx 1.005$, not $1.114$.
\end{center}

% ======================================================================
\section{Heat Kernel on Flat Torus (v70): Dead End}
\label{sec:heat_kernel}
% ======================================================================

\subsection{Setup}

The Seeley--DeWitt heat-kernel expansion on the $\psi$-circle (topologically
$T^1$ or more generally $T^3$ for the three imaginary directions) gives
subleading corrections to the effective action as a series in the background
curvature:
\begin{equation}
  \mathrm{Tr}\,e^{-t\Delta} \sim \frac{1}{(4\pi t)^{d/2}}
  \sum_{n=0}^\infty t^n\,c_n[g],
\end{equation}
where $c_n[g]$ are the Seeley--DeWitt coefficients.

\subsection{Key Computation}

On the flat torus $T^3$ (constant metric, zero curvature):
\begin{equation}
  R_{\mu\nu\rho\sigma} = 0
  \quad\Longrightarrow\quad
  c_n[g] = 0 \;\text{ for all } n \geq 1.
\end{equation}
Therefore all Seeley--DeWitt subleading corrections vanish identically on the
flat $\psi$-torus.  No correction of the form $R \approx 1 + c\alpha$ can arise
from the heat kernel on a flat background.

\subsection{Verdict}

\begin{center}
\textbf{Dead End} — flat $T^3$ has $c_n = 0$; no heat-kernel correction
generates $R \approx 1.114$.
\end{center}

% ======================================================================
\section{Summary: R Remains Open Hard Problem (v70)}
\label{sec:v70_summary}
% ======================================================================

After 25+ approaches across all versions, the $R$ factor remains an
\textbf{Open Hard Problem}.

\medskip
\begin{center}
\begin{tabular}{|l|l|l|}
\hline
\textbf{Approach (v70)} & \textbf{Result} & \textbf{Status} \\
\hline
Two-scale ratio $R_\psi^{(\alpha)}/R_\psi^{(\mathrm{EW})}$
  & $\beta \approx 163$ (no algebraic value) & Dead End \\
Modular correction near $\tau=i$
  & correction $\approx 1.005$ (factor 20 too small) & Dead End \\
Heat kernel on flat $T^3$
  & $c_n = 0$ (zero curvature) & Dead End \\
\hline
\end{tabular}
\end{center}

\medskip
\textbf{Best candidate (unchanged)}: $R = 1 + \alpha(N_{\mathrm{eff}} + \pi + 1/4)
\approx 1.1123$, error $0.15\%$.
\begin{itemize}
  \item $\pi$: candidate origin — $S^1$ angular integration.
  \item $1/4$: candidate origin — Weyl spinor (2 DOF / $2\pi^2$ from 4D).
  \item Neither has been derived from $\mathcal{S}[\Theta]$.
\end{itemize}

\textbf{Status}: $R$ is a \textbf{Motivated Conjecture (B3)} — numerical
observation only.  The two-loop beta-function approach has the correct
qualitative structure $R = 1 + c\alpha$; the coefficient $c$ is not
algebraically derived.

\textbf{Next step}: Explicit two-loop $\beta$-function coefficient on
$T^3 \times S^1$ (direct Feynman diagram computation in $\mathcal{S}[\Theta]$).

\medskip
\textbf{DERIVATION\_INDEX.md update (v70)}: R factor row — three additional
approaches (Two-Scale, Modular, Heat Kernel) all Dead End; best candidate
unchanged; status remains \textbf{Open Hard Problem}.

% ======================================================================
\section{Explicit Two-Loop $\beta$-Function Coefficient on $T^3 \times S^1$ (v71)}
\label{sec:twoloop_explicit}
% ======================================================================

% v71: Explicit computation of two-loop beta-function coefficient c.
% Approach: scalar field, SU(2), SU(2)×U(1) with N_f fermions.
% Result: no natural particle content reproduces c ≈ 15.39 — Dead End.

\subsection{Setup: Two-Loop Vacuum Diagram on $T^3$}

The UBT action on the $\psi$-circle $S^1$ at Kaluza-Klein level $n_R = 138$
is effectively a field theory on $T^3 \times S^1$ (three imaginary directions
times one compact real direction).  The two-loop beta function is controlled
by the two-loop vacuum bubble (sunset diagram):
\begin{equation}
The two-loop vacuum bubble (sunset diagram) contributes as:
\begin{equation}
  \Gamma^{(2)}\big|_{\text{2-loop}} \sim g^4 \!\!\int\!\! \frac{d^dp\,d^dk}{(2\pi)^{2d}}
  \frac{1}{p^2\,(k^2)\,((p+k)^2)},
  \label{eq:twoloop_integral}
\end{equation}
where $d = 3$ (the three compact directions of $\mathrm{Im}(\mathbb{H})$)
and the momenta are Kaluza-Klein sums $p_n = 2\pi n/R_\psi$.

The two-loop $\beta$-function coefficient $c_{2\text{loop}}$ is defined by:
\begin{equation}
  \beta_2\,\alpha^3 = -c_{2\text{loop}} \cdot \frac{\alpha^3}{(4\pi)^2},
  \quad
  R \;=\; 1 + \frac{\beta_2}{\beta_1}\alpha \;=\; 1 + c\alpha,
  \quad c = \frac{c_{2\text{loop}}}{c_{1\text{loop}}/(4\pi)^2}.
  \label{eq:c_2loop_def}
\end{equation}

\subsection{Case A: Scalar Field on $S^1$ (KK Sum)}

For a free scalar $\Theta$ on $S^1$ with KK modes summed over $n \in \mathbb{Z}$,
the standard two-loop result from the KK sum gives:
\begin{equation}
  c_{2\text{loop}}^{(\text{scalar})}
  = \frac{N_{\mathrm{eff}}}{2\pi}
  \approx \frac{12}{2\pi} \approx 1.91.
  \label{eq:c_scalar}
\end{equation}
\textbf{Check}: $c = 1.91 \neq 15.39$.
\begin{center}
  \textbf{Scalar field: $c \approx 1.91$ — Dead End.}
\end{center}

\subsection{Case B: Gauge Field $\mathrm{SU}(2)$ on $T^3$}

For a pure $\mathrm{SU}(2)$ gauge theory in $d=3$ dimensions, the two-loop
$\beta$-function coefficient involves the adjoint Casimir $C_A = 2$:
\begin{equation}
  c_{2\text{loop}}^{(\mathrm{SU}(2))}
  = \frac{11\,C_A}{3}
  = \frac{22}{3} \approx 7.33.
  \label{eq:c_su2}
\end{equation}
\textbf{Check}: $c = 7.33 \neq 15.39$.
\begin{center}
  \textbf{$\mathrm{SU}(2)$ gauge: $c \approx 7.33$ — Dead End.}
\end{center}

\subsection{Case C: $\mathrm{SU}(2) \times \mathrm{U}(1)$ with $N_f$ Fermions}

For a gauge theory with $\mathrm{SU}(2)$ gauge bosons and $N_f$ Dirac
fermion doublets, the two-loop coefficient is (using the standard result
$c_{\text{total}} = \frac{11}{3}C_A - \frac{4}{3}N_f T_F$, with $T_F = 1/2$
for the fundamental representation):
\begin{equation}
  c_{\text{total}} = \frac{11}{3}\cdot 2 - \frac{4}{3}\cdot N_f \cdot \frac{1}{2}
  = \frac{22}{3} - \frac{2}{3}N_f.
  \label{eq:c_total}
\end{equation}
Solving $c_{\text{total}} = 15.39$ for $N_f$:
\begin{equation}
  N_f = \frac{22/3 - 15.39}{2/3} = \frac{7.33 - 15.39}{0.667}
  \approx \frac{-8.06}{0.667} \approx -12.1.
  \label{eq:Nf_required}
\end{equation}
A \emph{negative} number of fermions is unphysical.

Alternatively, checking natural UBT particle content:
\begin{center}
\begin{tabular}{|l|l|l|}
\hline
\textbf{Content} & $N_f$ & $c_{\text{total}}$ \\
\hline
Pure gauge (no fermions) & 0 & $22/3 \approx 7.33$ \\
1 Weyl doublet & 1 & $22/3 - 2/3 = 20/3 \approx 6.67$ \\
3 generations, 2 doublets/gen & 6 & $22/3 - 4 = 22/3-12/3 = 10/3 \approx 3.33$ \\
3 gen, 4 Weyl doublets/gen (quarks+leptons) & 12 & $22/3 - 8 = 22/3-24/3 = -2/3 \approx -0.67$ \\
SM: 3 gen, 6 weak doublets total & $N_f = 6$ & $10/3 \approx 3.33$ \\
\hline
\end{tabular}
\end{center}
No entry reaches $c \approx 15.39$; all values are below $22/3 \approx 7.33$
and decrease with additional fermions.

\subsection{Case D: Including Higher-Spin KK Modes ($N_{\mathrm{eff}} = 12$)}

The full UBT spectrum has $N_{\mathrm{eff}} = 12$ effective bosonic modes.
Including all modes in the two-loop diagram (treating all 12 modes as
independent gauge bosons in an effective theory):
\begin{equation}
  c_{\text{full}} = \frac{11 \cdot N_{\mathrm{eff}}}{3}
  = \frac{11 \times 12}{3} = 44.
  \label{eq:c_full}
\end{equation}
This exceeds $15.39$ by a factor of $\sim 2.9$.  With mixed statistics
(fermion subtraction) it will decrease, not increase toward $15.39$.

\subsection{Summary: Two-Loop Diagram — Dead End}

\begin{center}
\begin{tabular}{|l|l|l|l|}
\hline
\textbf{Case} & \textbf{Content} & $c_{2\text{loop}}$ & \textbf{vs.\ target $15.39$} \\
\hline
A & Scalar on $S^1$ (KK) & $N_{\mathrm{eff}}/(2\pi) \approx 1.91$ & Too small ✗ \\
B & Pure $\mathrm{SU}(2)$ & $11C_A/3 \approx 7.33$ & Too small ✗ \\
C & $\mathrm{SU}(2)+N_f$ fermions & $7.33 - (2/3)N_f < 7.33$ & Decreases ✗ \\
D & All $N_{\mathrm{eff}}=12$ modes & $11\times12/3 = 44$ & Too large ✗ \\
\hline
Required & — & $c \approx 15.39$ & Target \\
\hline
\end{tabular}
\end{center}

\begin{center}
\textbf{Result (v71): Explicit two-loop $\beta$-function coefficient on $T^3 \times S^1$}\\
\textbf{does not reproduce $c \approx 15.39$ for any natural UBT particle content.}\\
\textbf{Dead End — standard two-loop gauge theory exhausted.}
\end{center}

\medskip
\textbf{Conclusion (v71)}: The two-loop $\beta$-function approach (B3) is
confirmed as a Dead End via explicit diagram computation:
\begin{itemize}
  \item Scalar: $c \approx 1.91$ (factor $\sim 8$ too small).
  \item Pure $\mathrm{SU}(2)$: $c \approx 7.33$ (factor $\sim 2$ too small).
  \item Adding fermions only decreases $c$; no physical fermion content
        reaches $15.39$.
  \item Including all $N_{\mathrm{eff}} = 12$ modes as gauge bosons overshoots
        ($c = 44$, factor $\sim 3$ too large).
\end{itemize}
The conjecture $c = N_{\mathrm{eff}} + \pi + 1/4 \approx 15.39$ does not
follow from any standard two-loop QFT computation on $T^3 \times S^1$.
This confirms Gap~B (R factor) remains \textbf{Open Hard Problem}.

\medskip
\textbf{DERIVATION\_INDEX.md update (v71)}: R factor row — explicit two-loop
diagram on $T^3 \times S^1$ computed for scalar (Case A), SU(2) gauge (Case B),
SU(2)+$N_f$ fermions (Case C), and all-mode gauge (Case D); none reproduces
$c \approx 15.39$; approach B3 confirmed Dead End.
Best candidate $R = 1 + \alpha(N_{\mathrm{eff}}+\pi+1/4)$ remains Numerical
Observation; status: \textbf{Open Hard Problem} (26+ approaches, all exhausted).

% ======================================================================
\section{What Would Close the Problem}
\label{sec:what_would_close}
% ======================================================================

\begin{enumerate}
  \item \textbf{Explicit two-loop diagram}: Compute the two-loop self-energy
    on the $\psi$-circle using the UBT action $\mathcal{S}[\Theta]$.
    Extract the coefficient $c$ algebraically.
  \item \textbf{Check}: Does $c = N_{\mathrm{eff}} + \pi + 1/4$ follow from
    the diagram?  If yes: Motivated Conjecture → Proved [L1].
  \item \textbf{If no}: Document the algebraic result and update Gap B status.
\end{enumerate}

% ======================================================================
\section{Summary and Status}
\label{sec:summary}
% ======================================================================

\begin{center}
\begin{tabular}{|l|l|l|}
\hline
\textbf{Result} & \textbf{Value} & \textbf{Status} \\
\hline
$R$ target (phenomenological) & $\approx 1.114$ & Empirical \\
Best algebraic candidate & $1 + \alpha(N_{\mathrm{eff}}+\pi+1/4) \approx 1.1123$ & Numerical Observation \\
Two-loop structure (qualitative) & Correct form $1 + c\alpha$ & B3 Motivated \\
Coefficient $c$ derived & $c = N_{\mathrm{eff}}+\pi+1/4$ & \OPEN\ Hard Problem \\
Two-loop diagram computed & — & \OPEN \\
\hline
\end{tabular}
\end{center}

\noindent\textbf{Conclusion}: The $R$ factor problem (Gap B) remains open.
The two-loop approach has correct qualitative structure ($R = 1 + c\alpha$),
but the coefficient $c$ has not been derived algebraically.
Status: \textbf{Motivated Conjecture (B3)} — $c = N_{\mathrm{eff}} + \pi + 1/4$
is the best numerical match but is not algebraically derived.

\medskip
\noindent\textbf{Update DERIVATION\_INDEX.md}: R factor row — B3 approach
documented; Numerical Observation for $c$ value; status remains
[OPEN HARD PROBLEM].

\end{document}
